\section{A reflection barrier at every binary word width}
\label{sec:prime-reflection}

The triple-relation family gives a direct linearization obstruction at its
specified cyclic dimensions and extension-compatible widths. The following
family has a different scope: every binary width $b\ge1$, odd prime side
length, and source rate tending to one half. It does not replace either the
above-half subcode construction or the triple relation-rank diagnostic.

Let $d=2m+1$ be an odd prime, and put
\[
 W_b=\{x\in(\mathbb F_2^b)^d:x_0=0,\ x_j=x_{-j}\ \forall j\},\qquad
 C_b=\{X:\text{every row of }X\text{ belongs to }W_b\}.
\]
All rotations act on entry columns and leave lane indices fixed.

\begin{theorem}[The reflection source]\label{thm:reflection-source}
\formal{CapacityReflectionSource.lean}{CapacityReflection.finrank_rowSpace,CapacityReflection.finrank_source,CapacityReflection.source_disjoint_views,CapacityReflection.source_shared}
The source $C_b$ has binary dimension $bdm$. Its distinct rotated views
$T^cC_b$, $T^eC_b$ intersect only in zero. It therefore admits one total
deterministic shared decoder for every receiver, and in particular the same
function can serve the two nonempty classes $\{0\}$ and
$\mathbb Z/d\setminus\{0\}$.
\end{theorem}

\begin{outline}
The $m$ pairs $\{j,-j\}$ of nonzero coordinates are free $b$-vectors, so
$\dim W_b=bm$ and $\dim C_b=dbm$. The formal equivalence chooses the
smaller positive representative of each pair and treats coordinate zero
separately.

A vector in $\sigma^cW_b$ is invariant under reflection
$j\mapsto-j-2c$ and vanishes at $-c$. Two such symmetries with distinct
centers give invariance under translation by $2(e-c)\ne0$. Because $d$ is
prime, this translation generates the whole cycle. Thus a vector in two
distinct translates is constant, and its zero coordinate forces it to be
zero. Apply this independently to every row and all its lanes. The formal
proof obtains any displacement as a multiple of the nonzero step in the
field $\mathbb Z/d$.

For a nonzero received view, disjointness makes its receiver phase and
underlying source matrix unique. Return that source's requested row.
Return zero at the zero view and choose arbitrary values off the union of
the views. This is one well-defined total function. It is not asserted
linear or efficiently computable.
\end{outline}

\begin{theorem}[Every-width linear subcode barrier]\label{thm:reflection-linear-barrier}
\formal{CapacityReflectionBarrier.lean}{CapacityReflection.one_class_finrank_le,CapacityReflection.two_class_finrank_le,CapacityReflection.two_class_codimension_lower,CapacityReflection.two_class_positive_deletion,CapacityReflection.deletion_rate_limit}
Let $C'\le C_b$ be any binary-linear subcode; it may couple every row and
every lane. If it has a single binary-linear shared decoder, then
\[
 \dim_{\mathbb F_2}C'\le b m(m+1).
\]
If instead it has two binary-linear shared decoders for the original
classes $\{0\}$ and $\mathbb Z/d\setminus\{0\}$, then
\begin{equation}\label{eq:reflection-two-class}
 \dim C'\le b\bigl(m(m+1)+d\bigr),\qquad
 \operatorname{codim}_{C_b}C'\ge b(m^2-d).
\end{equation}
The deletion bound is positive for odd primes $d\ge7$ and is asymptotic
to $bd^2/4$.
More generally, these linear-subcode bounds hold over every field of
characteristic two and at every odd side length $d=2m+1$. Primeness is
needed for the preceding deterministic-decoder construction.
\end{theorem}

\begin{outline}
First suppose there is one linear decoder $L$. Extend the coefficient
field to a splitting field $E$ for a primitive $d$th root of unity.
The formal construction takes the $E$-span of the embedded original basis
of $C'$ and proves that basis remains independent. The extended code
therefore has the same dimension. Extending the actual coordinate
coefficients of $L$ preserves its shared-decoder equations, and span
induction preserves the literal reflection and zero-coordinate equations
defining $C_b$. All row and lane coordinates are included. This step
extends only linear equations;
it makes no scalar-extension claim about the deterministic decoder of
\cref{thm:reflection-source}.

Since $d$ is odd, its image in the characteristic-two field is nonzero;
the formal primitive-character construction supplies the required
extension. Choose the Fourier basis $e_a(j)=\zeta^{aj}$,
$a\in\mathbb Z/d$. Its Vandermonde matrix is invertible. Apply this
\emph{same} change of basis to both physical indices, preserving the
transpose swap, and leave the lanes unchanged. Let $\mathcal F$ be the
one-index synthesis and define the column mode
$M_j(v)_{r,c,\ell}=\zeta^{jc}(\mathcal Fv)_{r,\ell}$.
Writing $X$ as the two-index synthesis of $Z$, its actual rotated view is
\[
 T^cX=\sum_j\zeta^{jc}M_j(\operatorname{col}_jZ).
\]
Apply $L$ and compare the independent characters in the receiver $c$
with the demanded row of $X$. This gives
$\mathcal F(\operatorname{row}_jZ)=L(M_j(\operatorname{col}_jZ))$.
Consequently, with the fixed linear maps
$B_j=\mathcal F^{-1}LM_j$, agreement takes the form
\begin{equation}\label{eq:reflection-frequency-equations}
 \operatorname{row}_j(Z)^{\mathsf T}
       =B_j\operatorname{col}_j(Z),\qquad j\in\mathbb Z/d,
\end{equation}
where $B_j:E^{db}\to E^{db}$ may mix every row and lane. No frequency
equation or dimension certificate is supplied as a premise: the formal
proof derives these equations from the original decoder.

In the Fourier basis a reflected row has equal coefficients at $a,-a$:
reflection exchanges those characters, and their independence forces
coefficient equality. Evaluating the row at physical coordinate zero sums
all coefficients. Equal nonzero pairs cancel in characteristic two, so
the zero-coordinate condition forces coefficient zero to vanish.
Conversely these coefficient conditions imply both original conditions.
Consequently
\[
 \operatorname{col}_0(Z)=0,\qquad
 \operatorname{col}_a(Z)=\operatorname{col}_{-a}(Z).
\]
The further Fourier change on the row index preserves these column
identities. Equation~\eqref{eq:reflection-frequency-equations} at zero now
forces $\operatorname{row}_0(Z)=0$ as well.

Choose representatives $1,\ldots,m$ of the nonzero frequency pairs.
Group the remaining rows and columns into those pairs. A $2$-by-$2$
block has duplicate columns and is specified by
$v_{a,b}\in E^{2b}$, its two entries in column $b$. The equations for
frequencies $a$ and $-a$ use the same input column, and show that the
whole compressed row $(v_{a,b})_b$ is a fixed linear function of the
compressed column $(v_{c,a})_c$. Equivalently, for fixed linear maps
$L_{a,b,c}:E^{2b}\to E^{2b}$,
\[
 v_{a,b}=\sum_c L_{a,b,c}(v_{c,a}).
\]
All cross-lane coefficients remain in these maps.

Supply the strict upper blocks $v_{a,b}$, $a<b$, and all diagonal blocks.
There are $\binom m2+m$ blocks, hence $bm(m+1)$ field coordinates.
Recover rows in descending order. In column $a$, the entries above the
diagonal were supplied, the diagonal was supplied, and those below it
belong to already recovered rows. The whole column is known and its fixed
linear maps determine row $a$. Thus the supplied-coordinate projection is
injective on the agreement space, proving its dimension bound. The formal
proof reuses the stronger descending block-recovery theorem, charging each
diagonal kernel by at most $2b$.

Finally suppose $C'$ has decoders $L_0,L_1$ for the two original classes.
In the \emph{original binary coordinates} impose
$L_1(X)=\operatorname{row}_0(X)$. This is one linear map to a $bd$-dimensional
row, so its kernel cuts $C'$ by at most $bd$ dimensions. On the resulting
subcode $C''$, the single decoder $L_1$ works also at receiver zero; it
already works at every nonzero receiver. The one-class bound gives
$\dim C'\le\dim C''+bd\le bm(m+1)+bd$. Subtract from
$\dim C_b=bdm$ and use $d=2m+1$ to get the stated deletion bound.
At $m\ge3$, $m^2-(2m+1)>0$, which covers the odd primes $d\ge7$.
\end{outline}

The source rate is $(d-1)/(2d)$, approaching one half from below, and the
statement covers every $b$ including $b=1$ and the canonical width $b(d)$.
Prime side length is used for the disjointness of rotated source views;
characteristic two is used for the Fourier zero coefficient. The
two-class partition in the linear bound is specifically singleton versus
complement. The linear bound itself needs only odd side length. The
result supplies neither a paid implementation of the nonlinear decoder
nor a full-cube or native impossibility theorem.
